\section{NeuroEvolution of Augmenting Topologies}

NeuroEvolution of Augmenting Topologies (NEAT), introduced by Stanley and Miikkulainen \cite{stanley2002evolving}, is an evolutionary algorithm that optimizes both the weights and the topology of artificial neural networks simultaneously. Traditional neuroevolution fixes the network architecture in advance and uses evolutionary or gradient-based methods only to adjust the weights. NEAT relaxes this constraint, allowing evolution to add neurons and connections incrementally, starting from a minimal network and growing complexity through mutation.

\subsection{Encoding and Historical Markings}

NEAT represents each network as a genome consisting of two gene lists: node genes and connection genes. A connection gene encodes a directed edge between two nodes, carrying a weight, an enable flag, and an \textit{innovation number} --- a global counter incremented each time a new structural mutation occurs anywhere in the population. Innovation numbers serve as historical markings that identify which structural mutations arose at the same point in evolutionary history, making it possible to align genomes of different topologies for crossover without requiring computationally expensive topological matching.

When two genomes mate, connection genes are aligned by their innovation numbers. Genes present in both parents (matching genes) are inherited randomly from either parent. Genes present in only one parent (disjoint or excess genes) are inherited from the more fit parent. This alignment procedure allows crossover between structurally different networks while preserving functional substructures.

\subsection{Structural Mutations}

NEAT introduces two structural mutations in addition to the standard weight perturbation:

\begin{itemize}
    \item \textit{Add connection}: a new connection is inserted between two previously unconnected nodes, with a randomly initialized weight and a new innovation number.
    \item \textit{Add node}: an existing connection is split by inserting a new node. The original connection is disabled, and two new connections are added --- one from the source node to the new node with weight 1, and one from the new node to the target with the original weight. This initialization preserves the network's existing behavior while providing a scaffold for further adaptation.
\end{itemize}

Because new structural mutations initially degrade fitness before they can be refined by evolution, NEAT protects innovation through \textit{speciation}: genomes are grouped into species based on a compatibility distance measure computed from the number of disjoint and excess genes and the average weight difference of matching genes. Fitness sharing within species ensures that novel structures are evaluated relative to topologically similar individuals rather than competing immediately against well-tuned incumbents.

\subsection{NEAT in Game-Playing Agents}

NEAT has demonstrated strong performance in sequential decision-making and game-playing tasks. Togelius and Lucas applied NEAT to evolve controllers for simulated car racing, demonstrating its capacity to discover reactive policies across continuous state spaces \cite{togelius2005evolving, togelius2006evolving}. More broadly, neuroevolution has been shown to match or exceed gradient-based methods on many game domains, particularly those with deceptive gradient landscapes or sparse rewards, where evolutionary exploration of topology provides an advantage \cite{risi2015neuroevolution}.

\section{CoDeepNEAT: Evolving Deep Architectures using NEAT}

The broader claim that neuroevolution can match human-designed architectures at competitive scale is supported by CoDeepNEAT (Coevolution DeepNEAT), introduced by Miikkulainen et al. \cite{miikkulainen2019evolving}. CoDeepNEAT extends NEAT to the problem of designing full deep learning architectures by introducing a two-level coevolutionary structure: a population of \textit{modules} (small subgraphs, each representing a repeatable network component such as a convolutional block) and a population of \textit{blueprints} (higher-level graphs whose nodes point to module species). During fitness evaluation, a blueprint is instantiated by replacing each of its nodes with a module drawn from the corresponding species, assembling a complete deep network that is then trained by gradient descent and scored by validation performance.

This bilevel coevolution mirrors the modular structure of successful hand-designed networks such as GoogLeNet and ResNet, which are built from repeated blocks, but discovers that structure automatically rather than by human intuition. Because small mutations to a module propagate to every blueprint node that references it, CoDeepNEAT can explore structurally diverse and deep architectures more efficiently than evolving a monolithic network would allow.

On the CIFAR-10 image classification benchmark, CoDeepNEAT achieved a test error of 7.3\%, comparable to the 6.4\% reported by the Bayesian hyperparameter optimization baseline of Snoek et al. Notably, the best evolved network trained significantly faster than the human-designed counterpart: it reached 20\% test error in 12 epochs versus more than 30 for the baseline, and converged around epoch 120 versus over 200. For language modeling on the Penn Treebank dataset, 25 generations of neuroevolution over LSTM connectivity improved test perplexity by 5\% over a vanilla LSTM by discovering a feedback skip connection between memory cells of two LSTM layers. This discover of innovative network architecture was possible only because the search was not constrained to a fixed topology, allowing evolution to explore structural mutations that gradient descent over a fixed architecture could not have found.

These results establish that NEAT-based evolutionary search over architecture space can discover both competitive performance and structurally novel solutions without expert intervention. For this thesis the key implication is one of scope: CoDeepNEAT operates at the scale of millions of parameters and full deep networks, yet the same underlying evolutionary mechanism applies equally to the small heuristic networks required for Sokoban A* search.

The most striking demonstration of the compactness potential of neuroevolution is Cuccu, Togelius, and Cudr\'{e}-Mauroux's \textit{Playing Atari with Six Neurons} \cite{cuccu2019playing}, which showed that a neuroevolved network of only 6--18 neurons, paired with a co-evolved compact state representation, matches or exceeds the performance of deep reinforcement learning agents using millions of parameters on several Atari benchmarks. Taken together, CoDeepNEAT and the six-neuron result benchmarks the range of what neuroevolution can achieve: at the large end, it matches expert-designed deep networks while at the small end, it discovers solutions orders of magnitude more compact than gradient-trained baselines. This motivates the central hypothesis of this thesis which is that NEAT can discover a heuristic network for Sokoban A* search that is drastically smaller than a fixed-topology baseline while remaining competitively accurate.
